\documentclass[12pt]{article}
\usepackage[letterpaper,margin=1.01in]{geometry}
\usepackage{amsmath,amssymb,amsfonts}
\usepackage{mathrsfs}
\usepackage{microtype}

\title{Homework: week 1}
\date{September 19, 2025}

\begin{document}
\maketitle

$R$: a commutative ring with a unit.

\bigskip
\noindent
\textbf{Problem 1.} Let $M$ be a $R$-module. Let $T^n(M) = M \otimes_R M \otimes_R \cdots \otimes_R M$ be the $n$-th tensor product. For $n = 0$, we put $T^0(M) = R$. Then
\[
T(M) = \bigoplus_{n\ge 0} T^n(M)
\]
is a (noncommutative) $R$-algebra, called the \emph{tensor algebra} of $M$. We define the \emph{symmetric
algebra}
\[
\operatorname{Sym}^\bullet(M) = \bigoplus_{n\ge 0} \operatorname{Sym}^n(M)
\]
to be the quotient quotient of the two-sided ideal generated by all expressions $x \otimes y - y \otimes x$,
for all $x,y \in M$; $\operatorname{Sym}^n(M)$ is called the $n$-th \emph{symmetric power}.

We also define the \emph{exterior algebra}
\[
\bigwedge(M) = \bigoplus_{n\ge 0} \bigwedge^n(M)
\]
to be the quotient of $T(M)$ by the two-sided ideal generated by all expressions $x \otimes x$ for
$x \in M$; $\bigwedge^n(M)$ is called the $n$-th \emph{Wedge power} of $M$. Now we assume that $F$ is a free
$R$-module of rank $r$. Then prove:

\begin{enumerate}
  \item $\bigwedge(M)$ is a skew commutative graded $R$-algebra, which means that if $u \in \bigwedge^k(M)$ and
  $v \in \bigwedge^l(M)$, then $u \wedge v = (-1)^{kl} v \wedge u$.
  \item $\operatorname{Sym}^k(F)$ and $\bigwedge^k(F)$ are also free. Find their ranks as $R$-modules.
  \item $\operatorname{Sym}^\bullet(F) \simeq R[\xi_1,\ldots,\xi_r]$.
  \item For every $k$, the multiplication map $\bigwedge^k(F) \otimes \bigwedge^{r-k}(F) \to \bigwedge^r(F)$ induces an isomorphism
  \[
  \bigwedge^k(F) \simeq (\bigwedge^{r-k}(F))^* \otimes \bigwedge^r(F),
  \]
  where $(\bigwedge^{r-k}(F))^* = \operatorname{Hom}_R(\bigwedge^{r-k}(F),R)$ is its dual module.
\end{enumerate}

\begin{enumerate}
  \setcounter{enumi}{4}
  \item Let $0 \to F' \to F \to F'' \to 0$ be a short exact sequences of free $R$-modules. Then for
  every $k$, there exists a decreasing finite filtration
  \[
  \operatorname{Sym}^k(F) = G^0 \supseteq G^1 \supseteq \cdots \supseteq G^k \supseteq G^{k+1} = 0
  \]
  such that $G^p/G^{p+1} \simeq \operatorname{Sym}^p(F') \otimes_R \operatorname{Sym}^{k-p}(F'')$.
  \item For every $k$, there exists a decreasing finite filtration
  \[
  \bigwedge^n(F) = G^0 \supseteq G^1 \supseteq \cdots \supseteq G^k \supseteq G^{k+1} = 0
  \]
  such that $G^p/G^{p+1} \simeq \bigwedge^p(F') \otimes_R \bigwedge^{k-p}(F'')$. In particular,
  \[
  \bigwedge^r(F) = \bigwedge^{r'}(F') \otimes_R \bigwedge^{r''}(F''),
  \]
  where $r'$ and $r''$ are ranks of $F'$ and $F''$ respectively.
\end{enumerate}

$A$: a ring with unit (possibly not commutative).

\bigskip
\noindent
\textbf{Problem 2.} Let $\Gamma_\bullet A$ be a $\mathbb{Z}_{\ge 0}$-filtration of $A$. Prove that if $\operatorname{gr}_\bullet A$
is noetherian, then so is $A$ itself.

\bigskip
Let $\Omega_\bullet A$ be a $\mathbb{Z}$-filtration of $A$. We define a topology on $A$ by declaring a subset to be
open precisely if it is a union of cosets $a + \Omega_i A$ for $a \in A$ and $i \in \mathbb{Z}$, called the filtration
topology of $A$.

\bigskip
\noindent
\textbf{Problem 3.} Prove that the filtration topology of $A$ is Hausdorff.

Hint: use the condition $\bigcap_i \Omega_i M = 0$.

\bigskip
A \emph{Cauchy sequence} in $A$ is a sequence $(m_n)_{n\ge 0}$ such that for every open set $U$ containing
$0$, there exists $c=c(U) \in \mathbb{N}$ such that if $i,j>c$ then $m_i-m_j \in U$. We say two Cauchy
sequences $(m_n)$ and $(m_n')$ are equivalent if for every open set $U$ containing $0$, there exists
$c=c(U) \in \mathbb{N}$ such that if $n>c$ then $m_n-m_n' \in U$.

Let $\widehat{A}$ be the set of all all equivalence classes of Cauchy sequences.

\bigskip
\noindent
\textbf{Problem 4.} Prove that $\widehat{A}$ is a ring and that we have a ring homomorphism
\[
A \to \widehat{A}, \quad a \mapsto [(a)],
\]
where $(a) = (a,a,\ldots,a,\ldots)$ is the constant sequence.

\bigskip
$A$ is called complete with respect to $\Omega_\bullet A$ if $A \simeq \widehat{A}$.

\bigskip
\noindent
\textbf{Problem 5.} Prove that if $A$ is complete and if $\operatorname{gr}_\bullet A$ is noetherian, then $A$ is also noetherian.

\bigskip
\noindent
\textbf{Problem 6.} Let $X$ be a smooth affine algebraic variety over $\mathbb{C}$. Prove
\[
\operatorname{gr}^F_\bullet D_X \simeq \operatorname{Sym}^\bullet T_X
\]
where $F_\bullet D_X$ is the order filtration and $T_X$ is the module of $\mathbb{C}$-derivations on $X$.
\end{document}
